\section{Conclusion}
\label{sec:conclusion}

Four questions were asked and four were answered.

\textbf{Is it hostile?} No. Six independent lines of static evidence, all
negative, including a byte-for-byte comparison against an independently verified
package. The most alarming artefact in the distribution---a utility that rewrites
the system \verb|hosts| file---turned out to be the strongest single piece of
evidence \emph{for} the package's authenticity.

\textbf{Is commercial anti-tamper present?} No, and it never was in this build.
The executable's Authenticode digest validates over its current bytes, so the
file is exactly what the publisher signed. Protection of that class lives inside
the executable and cannot be removed without invalidating the signature.
Three corroborating instruments agree, including a windowed entropy profile with
zero windows anywhere near the threshold.

\textbf{Can it be reconstructed without its installer?} Yes---in 37 minutes, to a
working 54\,GB tree, with every one of ten stages verified in advance against a
manifest the format publishes for free. The artefact that blocked the project for
two working sessions was a three-line advertisement fetcher containing no build
logic at all, and it was blocked not by a codec but by a missing configuration
file at a hard-coded path.

\textbf{Will it run on both GPUs?} Yes. The integrated GPU worked early. The
discrete GPU produced a failure with no error message, no crash, and no log
output, which turned out to be nine unpainted modal dialogs each reporting that
the vendor's shader compiler had segmentation-faulted. The cause was neither the
game, nor the driver, nor the hardware, nor the permissions gate that was
suspected three separate times. It was the build of the Direct3D~12 translation
layer---the one component chosen at the start of the project and never once
treated as a variable. Replacing it took one attempt.

The recurring lesson is not about graphics. It is that in a stack of a dozen
substitutable layers, the expensive mistake is rarely a wrong hypothesis; it is a
component that was never a hypothesis. Twice in this project---the installer
helper that re-extracted itself, and the translation layer that was assumed
fixed---the answer was in a place that had been silently excluded from the search
space at the outset.

Two smaller results are worth carrying away independently. A validated code
signature is a stronger, cheaper, and more honest way to establish the absence of
anti-tamper protection than any heuristic stack, and it should be the first test
rather than the last. And a pipeline stage that publishes its own preconditions
gives you a correctness check you neither had to design nor have to trust
yourself to have designed correctly.

\section*{Appendix: Reproduction}
\label{sec:appendix}

\subsection*{A. Establishing anti-tamper absence on an arbitrary binary}

\begin{enumerate}
\item Extract the certificate directory from the PE file, parse the PKCS\#7
      structure, and read the embedded digest.
\item Recompute the Authenticode hash over the file as it exists, excluding the
      checksum field, the certificate table entry, and the certificate data.
      Compare. If they match, stop---the question is answered.
\item Otherwise, profile entropy in windows of 1\,MiB over the code section.
      Report the window count above 7.5 and state the threshold used. Never
      report a whole-file figure alone.
\item Enumerate sections and the import table. A full, compiler-shaped import
      table is difficult for a protection layer to preserve.
\item Identify what protection \emph{is} present. ``No anti-tamper'' is not ``no
      DRM.''
\end{enumerate}

\subsection*{B. Recovering a compiled Inno Setup script}

Search for the 12-byte identifier \verb|rDlPtS|; read nine little-endian 32-bit
words; seek to \textbf{offset0}; skip the 64-byte banner; read CRC, stored size,
and compression flag; then read \verb|[crc32(4) + up to 4096 bytes]| chunks until
the stored size is satisfied. The concatenated result is raw LZMA1 whose first
five bytes are the properties block. Do \emph{not} follow the \verb|zlb\x1a|
magic near the start of the overlay; that is the file data.

\subsection*{C. Reassembling the distribution}

\begin{enumerate}
\item Copy the decompressor's configuration file to the root of the emulated
      system drive.
\item Patch the compression helpers to drop the \verb|Global\| namespace prefix,
      using a live patcher that polls for newly extracted copies.
\item Extract all six containers to their \emph{correct} destinations
      (Table~\ref{tab:dest}); two of the six do not go where the others do.
\item Read the worklists from the smallest container for the commands, and the
      carved installer script for the ordering and working directories.
\item Before each directory patch, read its manifest and record the expected
      input count, input byte total, and output size. Verify after.
\item Parallelise the per-file worklists; the cost is process startup, not
      compression.
\end{enumerate}

\subsection*{D. Bisecting a translation-layer failure}

\begin{enumerate}
\item Enumerate windows inside the compatibility layer before assuming a hang.
      A modal dialog that is never painted is indistinguishable from a deadlock
      by every other means.
\item Read assertion text precisely. An NTSTATUS is not an API return code, and
      the difference decides whether you are debugging a rejection or a crash.
\item Dump shaders on both vendors and compare hashes \emph{and} bytes. Identical
      hashes with differing bytes means the layer is emitting vendor-specific
      code and your cross-vendor control is not a control.
\item Vary the translation-layer build before exhausting its configuration
      surface, not after.
\item Clone the distribution with hardlinks and replace libraries by
      delete-then-copy. Never edit in place. Give each variant its own prefix,
      and copy the prefix's sibling metadata, not just the prefix.
\end{enumerate}

\subsection*{E. Catalogue of silent failures encountered}

Each of the following reported success, or reported nothing, while doing
something other than what was intended.

\begin{table}[htbp]
\caption{Silent failures encountered in this project}
\label{tab:silent}
\centering
\footnotesize
\begin{tabularx}{\columnwidth}{@{}L{3.1cm}X@{}}
\toprule
\textbf{Construct} & \textbf{What it actually does} \\
\midrule
\verb|grep -c| with a fallback & prints \verb|0| \emph{and} exits non-zero, so
  the fallback appends a second zero; all later comparisons error and read false \\
Process match on full name & \verb|comm| truncates to 15 characters; the pattern
  never matches and the process reads as dead \\
Substring match on command line & matches the searching shell itself \\
Static patch of installer helpers & bypassed; the installer re-extracts fresh
  copies at runtime \\
Prefix seeded from \verb|pfx/| only & missing sibling metadata aborts the
  launcher before the game starts \\
Library search under \verb|lib64| & the libraries are under \verb|lib|; the
  search returns a confident false negative \\
GUID searched in printed order & the first three fields byte-swap in memory \\
Emptiness test on a \verb|/sys| file & always false; sysfs files report zero size
  regardless of content \\
Graphics layer's adapter report & deliberately reports AMD while on NVIDIA \\
Launcher's own notification & claimed NVIDIA while running on the integrated GPU \\
EGL vendor pin & replaces the vendor list rather than extending it, disabling an
  unrelated API for an unrelated vendor \\
Setgid helper raising only the effective GID & the shell drops it at startup \\
Archive tool's size counter & 32-bit signed; goes negative past 2\,GB \\
Debugger window text & truncated to 14 characters; labels look corrupted \\
Compositor geometry versus screenshots & two different coordinate systems
  differing by the scale factor \\
\bottomrule
\end{tabularx}
\end{table}

\subsection*{F. Artefacts}

The harness scripts (launcher, single-configuration trial, distribution variant
builder, three configuration sweeps, and a shader bisector), the archive
extraction harness and its freestanding 32-bit build recipe, the installer-script
carver, the helper patcher, and the reassembly drivers for both data chunks are
retained alongside a running findings log. Defect and decision records are
maintained in the project's own issue and decision registers.
